% BANKS-SOURCE: pdf=113 print=103 kind=section
\subsection{Confinement of monopoles in the Higgs phase}

% BANKS-SOURCE: pdf=113 print=103 kind=prose
The duality, monopole, and four-index epsilon formulas in this section are the
$D=4$ specialization, with orientation $\epsilon^{0123}=1$. The lattice model
below is likewise four-dimensional and Euclidean.

Maxwell's equations, with both electric and magnetic current sources, have the
form
\begin{equation*}
  \partial^\mu F_{\mu\nu}=J_\nu,
  \qquad \partial^\mu\widetilde F_{\mu\nu}=\widetilde J_\nu,
\end{equation*}
and are invariant under the duality transformation
\begin{equation*}
  F_{\mu\nu}\longrightarrow\widetilde F_{\mu\nu}
  =\frac12\epsilon_{\mu\nu\lambda\kappa}F^{\lambda\kappa},
\end{equation*}
if we exchange the electric current, $J^\mu$, with the magnetic one,
$\widetilde J^\mu$.

To quantize the system we must make a choice. It turns out that consistency
requires the coupling strengths of the two types of currents to be inversely
proportional to each other. The only way we can have a controlled semi-classical
expansion is to insist

% BANKS-SOURCE: pdf=114 print=104 kind=prose
that the strongly coupled type of particle (monopoles) becomes very massive as
the coupling gets strong, so that virtual loops of these particles become
negligible despite the strong coupling. Indeed, naive calculations of radiative
corrections to the mass of these particles suggest that it is $o(1/g^2)$,
where $g$ is the weak coupling of the electrically charged particles. In fact,
as we shall see in Chapter 10, magnetic monopoles arise as classical soliton
configurations, with masses of order $1/g^2$, whenever the $U(1)$ of
electromagnetism is obtained by spontaneous breakdown of a simple group.
Virtual effects of these solitons are exponentially small in $g$, because of
form factors. We will generally use the terminology ``electric charges'' to
refer to those particles which are weakly coupled in a semi-classical expansion
and ``magnetic monopoles'' to refer to the strongly coupled objects.

The choice we make upon quantization is to solve the magnetic Maxwell equation
as
\begin{equation*}
  F_{\mu\nu}(x)=\frac12\epsilon_{\mu\nu\lambda\kappa}
  \int\dd^4y\,[f^\lambda(x-y)\widetilde J^\kappa(y)]
  +\partial_\nu A_\mu-\partial_\mu A_\nu,
\end{equation*}
where
\begin{equation*}
  \partial_\lambda f^\lambda=\delta^4(x-y).
\end{equation*}
We shall also insist that the support of $f^\lambda(x)$ is on a one-dimensional
curve connecting the origin to infinity. $f^\lambda$ is called the Dirac
string. We will generally assume that the curve is in fact a straight line.
Then we are constructing the magnetic monopole by viewing it as one end of an
infinitely long, infinitesimally thin solenoid. Dirac argued that, classically,
the solenoid is invisible, but in quantum mechanics (in anticipation of
Aharonov and Bohm), charged-particle paths that encircle the solenoid would
pick up a phase $e^{\ii e g}$ ($g$ is the magnetic charge), which would enable
one to locate the direction of the solenoid. If that were the case, a monopole
would not be a point particle. This conclusion can be avoided only if there is a
universal quantization rule for the coupling strengths of electric and magnetic
poles
\begin{equation*}
  eg=2\pi N.
\end{equation*}

Actually, if we consider the possibility that particles carry both electric and
magnetic charge (and follow Schwinger by calling such particles dyons), then the
general quantization rule is
\begin{equation*}
  q_1m_2-q_2m_1=2\pi N,
\end{equation*}
for some integer $N$. This follows because, as one particle encircles the
other's solenoid, the second particle encircles the first's solenoid in the
opposite direction. In this formula we have used $q$ and $m$ to represent the
electric and magnetic charges in units of $e$ and $1/e$, respectively, where
$e^2$ is the semi-classical expansion parameter for Maxwell's field.

We note in passing that the modern treatment of monopole physics uses the
theory of fiber bundles. Instead of introducing the singular Dirac string
directly, one introduces two vector potentials, each of which describes the
magnetic field of a monopole only over part of space. Each has a Dirac string
singularity, but only in the part of space where it is not used. The two regions
overlap, and, on the sphere at infinity, the overlap

% BANKS-SOURCE: pdf=115 print=105 kind=prose
consists of a small region around some equator of the sphere. The difference
between the two vector potentials is a pure gauge, but the gauge transformation
has a non-trivial winding around the equator,
\begin{equation*}
  \oint\dd x^\mu\,\partial_\mu\omega=2\pi m,
\end{equation*}
where $m$ is the magnetic charge.

We now quantize the theory by writing the Lagrangian
\begin{equation*}
  \mathcal L=-\frac{1}{4e^2}F^D_{\mu\nu}F^{D\mu\nu}
  +A_\mu J^\mu-\frac{1}{2\alpha}(\partial_\mu A^\mu)^2.
\end{equation*}
We can do the path integral over $A_\mu$, and, if we take the electric and
magnetic currents to be conserved, it is independent of the choice of gauge.
The result is (in Euclidean space)
\begin{equation*}
\begin{split}
Z[J,\widetilde J]={}&\exp\!\left[-\frac{e^2}{2}\int\dd^4x\,\dd^4y\,
  J^\mu(x)D(x-y)J_\mu(y)\right]\\
  &\times\exp\!\left[-\frac{2\pi^2}{e^2}\int\dd^4x\,\dd^4y\,
  \widetilde J^\mu(x)D(x-y)\widetilde J_\mu(y)\right]\\
  &\times\exp\!\left[\ii\int\dd^4x\,\dd^4y\,\dd^4z\,
  \epsilon_{\mu\nu\lambda\kappa}J^\mu(x)\partial^\nu D(x-y)
  f^\lambda(y-z)\widetilde J^\kappa(z)\right].
\end{split}
\end{equation*}
Here $f^\lambda=n^\lambda(n^\alpha\partial_\alpha)^{-1}$, where $n^\alpha$ is
the space-like direction of the Dirac vortex string attached to each monopole.
We have imposed the minimal Dirac quantization condition and the currents are
normalized to 1. When the monopole is stationary, at the origin, and we choose
$n^\lambda=\hat x^3$, then
\begin{equation*}
  \int\dd^4z\,\epsilon_{\mu\nu\lambda\kappa}f^\lambda(y-z)
  \widetilde J^\kappa(z)=\theta(y^3)\delta^2(\mathbf y)
  m\epsilon_{\mu\nu34}\delta(y^4-\tau).
\end{equation*}
We have chosen the parameter along the monopole world line to be its Euclidean
time coordinate. The phase is then non-vanishing for space-like charged-particle
trajectories, which encircle the vortex, and is equal to $e^{\ii qm}$. This is
the derivation of Dirac's quantization rule.

We conclude that a theory of light charged particles interacting with heavy
magnetic monopoles makes sense, at least in the approximation that the monopoles
are treated as singular classical sources. When we realize that monopoles are
forced on us, as classical solitons in non-abelian theories, we will understand
that this had to be the case. Those theories make it clear that a fully
relativistic theory of electric charges interacting with monopoles must be
consistent as well.

Now we come to an apparent conundrum. Let's consider the Higgs model, but add
in a coupling to a heavy magnetic monopole. Gauss' law tells us that there must
be a magnetic Coulomb field at infinity, but we have seen that there are no
massless fields in the theory. How do we get a long-range field? You might have
asked the same question about electric charge, but here the answer is easy. The
field $\phi$ has a vacuum expectation value,\footnote[9]{There is a subtlety here,
because this field is not gauge-invariant. Later on, when we discuss Wilson line
operators, we will see how to convert this argument into one about a
gauge-invariant variable.} but does not commute with the electric-charge
operator. The vacuum is therefore

% BANKS-SOURCE: pdf=116 print=106 kind=figure id=8.2
\begin{figure}[htbp]
  \centering
  % BANKS-SOURCE: pdf=116 print=106 kind=figure id=8.2
\begin{tikzpicture}[x=1cm,y=1cm,>=Stealth]
  \coordinate (L) at (-4.2,0);
  \coordinate (R) at (4.2,0);
  \shade[ball color=gray!28] (L) circle (0.48);
  \shade[ball color=gray!28] (R) circle (0.48);
  \node at (L) {\Large$+$};
  \node at (R) {\Large$-$};
  \foreach \y in {-1.05,-.78,-.51,-.24,0,.24,.51,.78,1.05} {
    \draw[->,line width=.55pt] (-3.75,\y) .. controls (-1.7,{\y+0.32}) and
      (1.7,{\y+0.32}) .. (3.75,\y);
  }
  \foreach \y in {-1.32,1.32} {
    \draw[->,line width=.55pt] (-3.8,\y) .. controls (-1.55,{\y+0.7}) and
      (1.55,{\y+0.7}) .. (3.8,\y);
  }
\end{tikzpicture}

  \caption{A flux tube between monopoles in a superconductor.}
  \label{fig:8.2}
\end{figure}

% BANKS-SOURCE: pdf=116 print=106 kind=prose
not an eigenstate of electric charge, but rather a superposition of different
charge states. The charge fluctuations in the vacuum screen any external source
of charge, so that it has no long-range field. There are no dynamical monopoles
in the theory, so this cannot be the explanation of the absence of a monopole
Coulomb field.

Instead, we must recognize that Gauss' law only relates the charge to the
integrated magnetic flux at infinity. In ordinary empty space, the lowest-energy
configuration satisfying this constraint is the rotationally invariant Coulomb
field. In the Higgs vacuum, the quanta of the magnetic field are massive, and it
turns out that it costs an energy per unit volume to force magnetic flux through
the system. The minimal energy configuration with a single point monopole is a
flux tube (Figure~\ref{fig:8.2}), which has energy per unit length.

To see this explicitly, let us take the monopole off to infinity, and search for
the minimal energy configuration with an infinite straight flux tube and no
source. The energy density for $x^3$-independent static configurations is
% BANKS-SOURCE: pdf=116 print=106 kind=equation id=8.1
\begin{equation}
  \mathcal E=\frac{1}{e^2}\left[\frac{F_{12}^2}{2}+|D_i\phi|^2
  +\frac{\lambda}{2}(|\phi|^2-\phi_0^2)^2\right].
  \label{eq:8.1}
\end{equation}

If $\lambda=1$, we can write this as
% BANKS-SOURCE: pdf=116 print=106 kind=equation id=8.2
\begin{equation}
\begin{split}
  \mathcal E=\frac{1}{e^2}\bigg[&\frac{(F_{12}+|\phi|^2-\phi_0^2)^2}{2}
  +|D_i\phi+\ii\epsilon_{ij}D_j\phi|^2\\
  &+\phi_0^2F_{12}-\ii\epsilon_{kl}\partial_k(\phi^*D_l\phi)\bigg].
\end{split}
  \label{eq:8.2}
\end{equation}
Note that the last two terms are total derivatives and contribute to the total
energy only through their values at infinity. We will find that $D_l\phi$ falls
off exponentially, so that the contribution of these terms to the energy is
just $\phi_0^2/e^2$ times the total magnetic flux through the plane. We can
minimize the energy by setting the two perfect squares in $\mathcal E$ to zero,
so that this flux contribution is the total energy of the vortex. The

% BANKS-SOURCE: pdf=117 print=107 kind=prose
minimum energy solutions are determined by the Bogomolnyi--Prasad--Sommerfield
(BPS) \cite{banks-ref-82,banks-ref-83} equations\footnote[10]{For $\lambda=1$, the Higgs Lagrangian
is the bosonic sector of a supersymmetric model, and the BPS equations are the
conditions that the vortex solution preserve some of the supersymmetry.}
% BANKS-SOURCE: pdf=117 print=107 kind=equation id=8.3
\begin{equation}
  (D_i+\ii\epsilon_{ij}D_j)\phi=0,
  \label{eq:8.3}
\end{equation}
% BANKS-SOURCE: pdf=117 print=107 kind=equation id=8.4
\begin{equation}
  F_{12}=\phi_0^2-\phi^2.
  \label{eq:8.4}
\end{equation}
The solution with flux $n$ is
\begin{equation*}
  \phi=e^{\ii n\theta}f(r),
  \qquad A_1+\ii A_2=-\ii e^{\ii\theta}\frac{a(r)-n}{r},
\end{equation*}
where
\begin{equation*}
  f'=\frac{a}{r}f,
  \qquad a'=r(f^2-\phi_0^2).
\end{equation*}
The boundary conditions are
\begin{equation*}
  a\to0,\quad f\to\phi_0,\qquad\text{for }r\to\infty;
\end{equation*}
\begin{equation*}
  a\to n+O(r^2),\quad f\to r^n(1+O(r^2)),
  \qquad\text{for }r\to0.
\end{equation*}
It is easy to see, by linearizing the equations at infinity, that $f$ approaches
its asymptotic value exponentially. According to the second BPS equation, this
means that the magnetic flux is confined to a small region around $r=0$.
$\phi_0$ sets the scale for this region.

\BanksImplicitHook{I-CH08-006}

The qualitative nature of the solution is similar for other values of $\lambda$.
The Nielsen--Olesen vortex \cite{banks-ref-84} that we have constructed is the global
minimum-energy configuration with fixed flux. The fact that magnetic flux
penetrates the Higgs vacuum in thin tubes was first discovered, both
theoretically and experimentally, in the study of superconductivity. There the
expulsion of magnetic field from the superconducting ground state is called the
Meissner effect, and the flux-tube solution was found by Abrikosov \cite{banks-ref-85}.

The significance of the Meissner effect for the physics of dynamical monopoles
in the Higgs vacuum is profound. Isolated monopoles will have infinite energy.
Monopole--anti-monopole pairs will experience a linear confining potential
(energy per unit length for a flux tube connecting them) at large separation.
They will be permanently bound into states with a characteristic size determined
by the energy density of the flux tube.

At the time the Nielsen--Olesen vortex was discovered, experimental physics had
revealed a puzzle about the nature of quarks. Spectroscopy had given strong
indications of a quark substructure underlying hadrons. High-energy scattering
experiments, both deep inelastic scattering of leptons and nucleons and
electron--positron annihilation into hadrons, had given dramatic confirmation to
the idea that hadrons are bound states of quarks and that quarks behave almost
like free particles when they are close

% BANKS-SOURCE: pdf=118 print=108 kind=prose
together (asymptotic freedom). But quarks had never been seen as free particles,
despite the fact that experiments were being done at energies above the
characteristic scale of the strong interactions.

By 1973, the apparently free behavior of quarks in short-distance experiments
had been understood in terms of quantum chromodynamics (QCD), the $SU(3)$ color
gauge theory of strong interactions. As we will learn in Chapter 9, the
effective gauge coupling is scale-dependent and goes to zero at short distance.
Conversely, as the separation between quarks gets larger, it grows, up until a
length scale of order $(150\,\mathrm{MeV})^{-1}$, whereupon the perturbative
calculations on which this observation is based break down. Many people
speculated that the coupling became infinitely strong and led to the
confinement of quarks.

Actual calculations showing the confinement of quarks were first done (in four
dimensions) by Wilson, in a latticized form of the gauge theory \cite{banks-ref-86}. Shortly
thereafter, 't Hooft \cite{banks-ref-87} and Mandelstam \cite{banks-ref-88} noted that one could obtain an
intuitive picture of confinement by treating the Cartan subgroup of $SU(3)$ as
electromagnetism (actually two different $U(1)$ gauge fields since the group has
rank two), and postulating that the vacuum state of the theory was in a magnetic
Higgs phase. Using electric--magnetic duality and the calculations we have done
for monopoles in the Higgs phase, we see that such a state would lead to
confinement of quarks by electric flux tubes.

Much work in lattice gauge theory has gone into trying to validate the
't Hooft--Mandelstam picture of confinement. The results are a trifle
ambiguous. Nonetheless, this is a good qualitative picture of what is going on
in QCD. The picture has been verified both in lattice models and in partially
soluble supersymmetric gauge theories. We will now examine a simple lattice
model, which will introduce the reader to the techniques of lattice gauge theory
and exhibit the 't Hooft--Mandelstam mechanism for quark confinement in a
striking manner \cite{banks-ref-89,banks-ref-90,banks-ref-91,banks-ref-92}.

In a lattice theory (we will work in four-dimensional Euclidean space), the
functional integral is replaced by an integral over variables defined on a
hypercubic (for simplicity) lattice, and derivatives are replaced by finite
differences. In a lattice gauge theory, the fundamental variable is the parallel transporter, or Wilson line
between nearest-neighbor points on a lattice. We will denote this by
$U_\mu(x)$, where $x$ is a lattice point and $\mu$ one of the eight independent
directions on the lattice. It connects the points $x$ and $x+\hat\mu$. Under a
gauge transformation,
\begin{equation*}
  U_\mu(x)\longrightarrow V(x)U_\mu(x)V^\dagger(x+\hat\mu).
\end{equation*}
We also insist that $[U_\mu(x)]^{-1}=U_{-\mu}(x+\hat\mu)$. We take these
variables to be matrices, which are group elements in the fundamental
representation (we will deal only with $U(N)$ groups in this text).

Given gauge-variant fields $\phi(x)$, transforming in the representation $R$ of
the gauge group, we can write gauge-covariant finite differences
\begin{equation*}
  \phi(x)-U_\mu^{(R)}(x)\phi(x+\hat\mu),
\end{equation*}
where the matrix that appears is the $R$ representation of the group element
$U_\mu(x)$. These can be combined to form gauge-invariant actions. Gauge-
invariant actions constructed

% BANKS-SOURCE: pdf=119 print=109 kind=prose
solely from the gauge fields are functions of traces of Wilson loops. The
simplest such loop is the plaquette in the $\mu\nu$ plane,
\begin{equation*}
  U_{\mu\nu}(x)=U_\mu(x)U_\nu(x+\hat\mu)U_{-\mu}(x+\hat\mu+\hat\nu)
  U_{-\nu}(x+\hat\nu).
\end{equation*}
Since it goes around a closed curve, its trace is gauge-invariant.

All of this simplifies for abelian gauge groups, for which we can write
$U_\mu(x)=e^{\ii\theta_\mu(x)}$ and treat $\theta_\mu$ much as we would a gauge
potential in the continuum. We will consider the group $\mathbb Z_N$.

To couple a lattice gauge theory to a heavy external particle, we let $x_\mu(s)$
be the Euclidean path that the particle follows. Here $s$ is a discrete
parameter that counts the steps in the path. In order to define a gauge-invariant
quantity, we insist that the path be closed. Physically this corresponds to
studying a heavy charged particle by creating a particle--anti-particle pair,
separating them, and bringing them back together to annihilate. The observable
which computes the gauge-theory contribution to the action of this path is just
the corresponding Wilson loop $W[x(s)]$. The gauge charge of the heavy particle
is encoded in the choice of representation for the $U_\mu$ matrices in the
Wilson loop.

To compute the potential energy of the pair a distance $R$ apart we choose
(Figure~\ref{fig:8.3}) a curve containing long straight segments along which
only $x_0(s)$ changes, and the distance between the two straight segments is
$R$.
% BANKS-SOURCE: pdf=119 print=109 kind=figure id=8.3
\begin{figure}[htbp]
  \centering
  % BANKS-SOURCE: pdf=119 print=109 kind=figure id=8.3
\begin{tikzpicture}[x=1cm,y=1cm,>=Stealth]
  \draw[->,line width=.7pt] (-2.0,-2.8) -- (2.0,-2.8);
  \draw[->,line width=.7pt] (2.0,-2.8) -- (2.0,2.8);
  \draw[->,line width=.7pt] (2.0,2.8) -- (-2.0,2.8);
  \draw[->,line width=.7pt] (-2.0,2.8) -- (-2.0,-2.8);
  \node[below left] at (-2.0,-2.8) {$0$};
  \node[below right] at (2.0,-2.8) {$R$};
  \node[right] at (2.0,0) {$T$};
\end{tikzpicture}

  \caption{The Wilson loop for the static potential.}
  \label{fig:8.3}
\end{figure}
If $T$ is the length of the long segments and we have a linear confining
potential above some distance $R_c<R$, then we expect
\begin{equation*}
  \langle W[x(s)]\rangle\longrightarrow e^{-\kappa RT},
\end{equation*}
as $T\to\infty$. This is a special case of a more general and more geometric
formula
\begin{equation*}
  \langle W[x(s)]\rangle\longrightarrow e^{-\mathcal A},
\end{equation*}
% BANKS-SOURCE: pdf=120 print=110 kind=prose
where $\mathcal A$ is the area of the minimal surface spanning the loop.
Confinement is equivalent to Wilson's area law for large-area Wilson loops.

It is worth noting that, in non-abelian gauge theories, we do not expect to see
confinement for heavy particles in generic representations of the gauge group.
Indeed, such theories contain dynamical gauge bosons, which might be able to
combine with the heavy particle to form singlet bound states. These singlets can
propagate far away from each other without feeling a confining force. However,
if the group is $SU(N)$ it has a non-trivial center $\mathbb Z_N$. The adjoint
representation does not transform under the center, so it cannot screen
particles in representations that do transform.\footnote[11]{If the theory
contains dynamical particles in the fundamental representation, then any heavy
source can be screened by creation of pairs of these particles. This is what
happens for quarks in the real world.}

Since the essence of confinement in an $SU(N)$ theory is associated with the
$\mathbb Z_N$ center, it seems worthwhile to study the pure $\mathbb Z_N$
lattice gauge theory, and that is what we will do. The variables of this theory,
$\theta_\mu(x)$, are integers modulo $N$, one for each link on the lattice. A
group element is
\begin{equation*}
  U_\mu(x)=\exp\!\left(\frac{2\pi\ii\theta_\mu(x)}{N}\right).
\end{equation*}
We will choose a partition function of the form
\begin{equation*}
  Z=\int\prod_{x,\mu}[\dd\theta_\mu(x)]
  \prod_{x,\mu\nu}f(\theta_{\mu\nu}(x)).
\end{equation*}
The ``measure'' $[\dd\theta]$ is just the instruction to sum each link variable
over its $N$ possible values. The gauge-invariant field strengths
$\theta_{\mu\nu}$ are defined by
\begin{equation*}
  \theta_{\mu\nu}(x)=\theta_\mu(x+\hat\nu)-\theta_\mu(x)
  -\theta_\nu(x+\hat\mu)+\theta_\nu(x)
  \equiv\Delta_\nu\theta_\mu-\Delta_\mu\theta_\nu.
\end{equation*}
Like the original link variables $\theta_\mu$, these add modulo $N$.

We can write the most general function $f$ of these $\mathbb Z_N$-valued
variables as a finite Fourier series
\begin{equation*}
  f(\theta)=\frac1N\sum_l e^{\ii l\theta}g(l),
\end{equation*}
where the variable $l$ is also $\mathbb Z_N$-valued. If we insert the Fourier
transform for each $\theta_{\mu\nu}$, and integrate by parts, we can do the
$\theta_\mu$ ``integrals'' exactly, obtaining
\begin{equation*}
  Z=\int[\dd l_{\mu\nu}(x)]\prod_x[g(l_{\mu\nu}(x))]
  \delta[\Delta_\nu l_{\mu\nu}(x)],
\end{equation*}
where we have used the Einstein summation convention for the $\nu$ index inside
the delta function. The ``functional delta function,'' $\delta[\,]$, is the
product over all links, $(x,x+\hat\mu)$, in the lattice, of the
$\mathbb Z_N$-valued Kronecker delta of the variables
$\Delta_\nu l_{\mu\nu}(x)$. The Wilson-loop expectation value is the ratio of
two such integrals over $\theta_\mu(x)$. The argument of the functional delta
function in the numerator is $\Delta_\nu l_{\mu\nu}-qw_\mu(x)$. $q$ is the
$\mathbb Z_N$ charge of the static particle, and $w_\mu(x)$ is 1 if the Wilson
loop goes through the link $(x,x+\hat\mu)$ in the positive direction, $-1$ if
it goes through in the negative direction, and 0 if the Wilson loop does not
pass through that link.

% BANKS-SOURCE: pdf=121 print=111 kind=prose
To convert this to a form making its relation to electrodynamics clear, we
introduce redundant variables, to write the sums over $\mathbb Z_N$ variables
in terms of ordinary integer sums. Thus, we write
$l_{\mu\nu}\to l_{\mu\nu}+Nq_{\mu\nu}$, where $l$ and $q$ now take on arbitrary
integer values. Shifting $l$ by a multiple of $N$ can be compensated by a shift
of $q$, so, if all functions are just functions of this combination, the double
sum over integers just produces an infinite number of copies of the sum over
$\mathbb Z_N$ variables. This infinity cancels out in the ratios that define
expectation values. In order to preserve the $\mathbb Z_N$ character of the
delta functional, we introduce an integer-valued link variable $e_\mu(x)$ and
write the functional as
\begin{equation*}
  \delta[\Delta_\nu l_{\mu\nu}(x)-Ne_\mu(x)].
\end{equation*}
The ``current'' $e_\mu$ must be conserved in order to satisfy this constraint.

For convenience, we will choose the weighting function
$g(l_{\mu\nu})=e^{-g^2l_{\mu\nu}^2}$. Other choices give similar results. When
$g^2$ is very large, non-zero values of $l_{\mu\nu}$ are suppressed. If we
introduce a Wilson loop with $\mathbb Z_N$ charge 1, this shifts the delta
function in the numerator functional integral to
$\delta[\Delta_\nu l_{\mu\nu}-w_\mu]$. $w_\mu=\pm q$ or zero. The delta-function
constraint in the numerator now forces $l_{\mu\nu}$ to be $q$ over some area
bounded by the Wilson loop. For large $g^2$ we obviously want to pick the
minimal area. In the denominator we simply set $l_{\mu\nu}=0$ for large $g^2$.
We find that the Wilson loop falls off like $e^{-g^2\mathcal A}$, where
$\mathcal A$ is the minimal area it bounds.

\BanksImplicitHook{I-CH08-007}

In order to understand what happens for smaller values of $g^2$, we introduce
the Poisson summation formula:
\begin{equation*}
  \sum_l f(l)=\sum_m\int\dd a\,f(a)e^{2\pi\ii ma}.
\end{equation*}
This formula is true because the sum over $m$ gives rise to a sum of delta
functions of $a$ concentrated on the integers $a=l$. If
$f(a)=e^{-g^2a^2}$, we can do the Gaussian integral and obtain
\begin{equation*}
  \sum_l e^{-g^2l^2}=\sqrt{\frac{\pi}{g^2}}\sum_m
  e^{-\pi^2m^2/g^2}.
\end{equation*}

We will use this formula after solving the constraint equation for $l_{\mu\nu}$
as
\begin{equation*}
  l_{\mu\nu}=\bigl(n_\nu(n\cdot\Delta)^{-1}j_\mu
  -n_\mu(n\cdot\Delta)^{-1}j_\nu\bigr)
  +\epsilon_{\mu\nu\lambda\kappa}\Delta^\lambda l^\kappa.
\end{equation*}
$l_\mu$ is an integer-valued link variable, with an obvious gauge ambiguity.
$j_\mu=Ne_\mu$ in the denominator functional integral and
$j_\mu=Ne_\mu+qw_\mu$ in the numerator. We want to apply the Poisson summation
formula to the sums over $l_\mu(x)$. In principle we must eliminate the gauge
ambiguity first, and do so (for example by picking an axial gauge like
$l_3=0$) in a way that preserves the integer nature of the variable. However,
if we introduce the Poisson variables $m_\mu$ by a factor
\begin{equation*}
  \exp\!\left(2\pi\ii\sum_x l_\mu m_\mu\right),
\end{equation*}
and insist that $\Delta_\mu m_\mu=0$, then the answers will be independent of
the choice of gauge. The result is a lattice action depending on three types of
variable, $e_\mu$, $m_\mu$, and $a_\mu$.

% BANKS-SOURCE: pdf=122 print=112 kind=prose
The former are integer-valued conserved currents on the lattice, while the
latter is a lattice version of Maxwell's field. We have
\begin{equation*}
  S=\sum_x\left[g^2\left(\Delta_\mu a_\nu-\Delta_\nu a_\mu
  -\epsilon_{\mu\nu\lambda\kappa}n^\lambda(n\cdot\Delta)^{-1}j^\kappa\right)^2
  +2\pi\ii m_\mu a_\mu\right].
\end{equation*}
This is a discrete analog of the action for electric and magnetic monopoles
interacting via the Maxwell field. If we are interested in the small-$g^2$
limit, we should carry out a duality transformation and describe $j_\mu$ as an
electric current and $m_\mu$ as a magnetic current. The resulting action is
\begin{equation*}
  S=\sum_x\left[\frac{1}{4g^2}
  \left(\Delta_\mu A_\nu-\Delta_\nu A_\mu
  -\epsilon_{\mu\nu\lambda\kappa}n^\lambda(n\cdot\Delta)^{-1}2\pi m^\kappa\right)^2
  +\ii j_\mu A_\mu\right].
\end{equation*}
If we do the path integral over $A$, we get a non-local action for the currents
\begin{equation*}
\begin{split}
  S_{\mathrm{NL}}={}&\frac{g^2}{2}\sum_{x,y}[w_\mu+Ne_\mu](x)D(x-y)
  [w_\mu+Ne_\mu](y)\\
  &+\frac{2\pi^2}{g^2}\sum_{x,y}m_\mu(x)D(x-y)m_\mu(y)+\ii\Theta_D.
\end{split}
\end{equation*}
where $\Theta_D$ is the Dirac phase factor we have discussed above in the
continuum. $D(x-y)$ is the Green function of the four-dimensional lattice
Laplacian.

We treat this expression in the following way: keep the term with $x=y$ in
$D$ as is, and rewrite the rest as a path integral over $A$ with a modified
kinetic term, whose inverse is $D(x-y)-D(0)$. The sum over $e_\mu$ and $m_\mu$
is now weighted by
\begin{equation*}
  \exp\!\left[-N^2g^2D(0)\sum_xe_\mu(x)^2
  -\frac{4\pi^2}{g^2}D(0)\sum_xm_\mu(x)^2\right].
\end{equation*}
For small $g^2\ll1/N^2$, the non-zero values of $m_\mu$ are highly suppressed,
while the sum over $e_\mu$ is free. We do it by introducing a Lagrange
multiplier to enforce the constraint $\Delta_\mu e_\mu=0$, via a term
\begin{equation*}
  \exp\!\left(\ii\sum_x\theta(\Delta_\mu e_\mu)\right).
\end{equation*}
The sum over $e_\mu$ gives an action of the form
\begin{equation*}
  S=\sum_x\left[F(\Delta_\mu\theta-NA_\mu)+A_\mu w_\mu
  +\frac{1}{4g^2}\left(A_{\mu\nu}-F^D_{\mu\nu}(m)\right)^2
  +\frac{4\pi^2}{g^2}m_\mu^2\right].
\end{equation*}
$A_{\mu\nu}$ is the lattice field strength and $F^D_{\mu\nu}(m)$ is the Dirac
string field corresponding to the magnetic current $m_\mu$. This looks like a
Higgs model coupled to magnetic monopoles, and for small $g^2$ we can perform
the path integral semi-classically. The result is that the magnetic-monopole
loops are confined to small size and are rare. The Wilson loop just has a
perimeter law.

On the other hand, for large $g^2$ we instead apply this treatment to the
$m_\mu$ loops, getting a magnetic Higgs phase. Electric loops of $e_\mu$ are
generally suppressed. However, if the Wilson loop $w_\mu$ bounds a large area,
and if it has charge $q$ (mod $N$) rather than charge 1, then it pays to cancel
it to the smaller of $q$ and $N-q$. We get an area law

% BANKS-SOURCE: pdf=123 print=113 kind=prose
with a coefficient that is largest for $q\sim N/2$. We have already analyzed
the large-$g^2$ regime in terms of the $l_{\mu\nu}$ variables. Now we realize
that we were just describing the magnetic Higgs phenomenon. $l_{\mu\nu}$ is
just the quantized electric flux which must pay an energy price per unit length
to penetrate the magnetic Higgs vacuum.

Finally, when $N$ is large, one can argue that there is a regime of $g^2$ where
both the electric and magnetic currents, $e_\mu$ and $m_\mu$, are suppressed,
and the theory behaves like free electrodynamics. This Coulomb phase exists
roughly when
\begin{equation*}
  4\pi^2c_1>g^2>\frac{c_2}{N^2},
\end{equation*}
where the $c_i$ are constants of order 1. For large $N$, the $\mathbb Z_N$
lattice gauge model exhibits the main phases of gauge theories and the
electric--magnetic duality between Higgs and confinement phases. Other possible
phases exist in more complicated non-abelian theories, in which the particle
whose field exhibits the Higgs mechanism at strong coupling has both electric
and magnetic charge at weak coupling. These are called oblique confinement
phases. Non-abelian theories have also been shown to exhibit conformally
invariant phases that are not free-field theories. These go under the name of
non-abelian Coulomb phases.
